\documentclass[10pt]{article}
\usepackage[margin=0.84in]{geometry}
\usepackage[T1]{fontenc}
\usepackage{amsmath,amssymb,amsthm}
\usepackage{graphicx}
\usepackage{booktabs}
\usepackage[colorlinks=true,linkcolor=blue,citecolor=blue,urlcolor=blue]{hyperref}

\newtheorem{theorem}{Theorem}
\newtheorem{lemma}{Lemma}
\newcommand{\E}{\mathbb E}
\newcommand{\R}{\mathbb R}
\newcommand{\Z}{\mathbb Z}
\newcommand{\one}{\mathbf 1}

\title{Off-diagonal singularity of the current--current\\
correlation measure: two explicit ergodic models}
\author{Run \texttt{fa\_banach\_001} --- \texttt{agent\_lane\_11}}
\date{13 August 2026}

\begin{document}
\maketitle

\noindent\textbf{Status.} Candidate full counterexample to the broad
formulation; likely valid subject to human review. A second, completely
localized model shows that localization plus a Wegner estimate is still
insufficient. Neither result settles the narrower i.i.d.-per-physical-site
Anderson subclass.

\section{Source question}

Combes, Germinet, and Hislop~\cite{cgh} define the positive current--current
correlation (ccc) measure by
\begin{equation}\label{eq:source-ccc}
 M(\Delta_1,\Delta_2)
 =\mathcal T\bigl(vE_H(\Delta_1)vE_H(\Delta_2)\bigr),
 \qquad v=i[H,X],
\end{equation}
and ask on page~5 whether a density $m(E_1,E_2)$ exists when $E_1\ne E_2$.
They note the positive theorem of Bellissard--Hislop under strong disorder and
an analytic single-site density~\cite{bh}.

\begin{center}
\includegraphics[width=0.96\textwidth]{figures/open_problem_crop.png}
\end{center}

\noindent The unrestricted answer is negative for a scalar model with standard
one-site covariance. A complementary unit-cell model retains complete
localization, a bounded-density Wegner estimate, deterministic finite-range
hopping, and a randomness-independent velocity.

\section{Proof intuition}

The first mechanism is Floquet-theoretic. A random choice between the two
phases of a period-two potential is an ergodic family under ordinary lattice
translations. Interband current transitions conserve quasi-momentum, so their
two energy coordinates trace a curve, rather than filling the energy plane.

The second mechanism shows why localization does not repair the issue. Put one
random scalar $\omega_n$ on a two-level cell. It shifts both cell energies
together, while a fixed hopping keeps their gap equal to $2t$. The velocity
exchanges the two eigenvectors, so averaging spreads ccc mass along two lines.
In both cases a square of side $\varepsilon$ captures order $\varepsilon$ mass
rather than order $\varepsilon^2$ mass.

\section{An exact scalar counterexample with standard translations}

Let $a,t,s>0$ satisfy $2s>t^2/a$. On scalar $\ell^2(\Z)$ define, for
$q\in\{0,1\}$,
\begin{equation}\label{eq:periodic-model}
 (K_q\psi)(j)=s\bigl(\psi(j+2)+\psi(j-2)\bigr)
 +t\bigl(\psi(j+1)+\psi(j-1)\bigr)+a(-1)^{j+q}\psi(j).
\end{equation}
Give the two values of $q$ equal probability. Standard translation by one
site sends $K_q$ to $K_{q+1\bmod2}$, so this two-point family is covariant.
The flip action is transitive and hence ergodic. Since only the diagonal
potential changes with $q$, $i[K_q,X]$ is independent of $q$.

\begin{theorem}\label{thm:periodic}
The ccc measure of \eqref{eq:periodic-model} is not absolutely continuous on
the off-diagonal. More strongly, at a suitable off-diagonal point
$(E_1^0,E_2^0)$,
\[
 \frac{M([E_1^0,E_1^0+\varepsilon],
             [E_2^0,E_2^0+\varepsilon])}{\varepsilon^2}
 \longrightarrow+\infty.
\]
Thus the unrestricted off-diagonal density assertion is false on scalar
$\ell^2(\Z)$ with the standard translation representation.
\end{theorem}

\begin{proof}
It is enough to use the phase $q=0$, since the two phases are translates. In a
two-site Floquet decomposition, with $k\in[-\pi,\pi]$,
\[
 K(k)=2s\cos k\,I+d(k)\cdot\sigma,
 \qquad d(k)=\bigl(t(1+\cos k),t\sin k,a\bigr).
\]
Put
\[
 R(k)=\sqrt{a^2+4t^2\cos^2(k/2)},\qquad
 \lambda_\pm(k)=2s\cos k\pm R(k).
\]
In the same representation the velocity is
\[
 v(k)=-4s\sin k\,I+g(k)\cdot\sigma,
 \qquad g(k)=\bigl(-t\sin k,-t(1-\cos k),0\bigr).
\]
The scalar part has zero interband matrix element. For a two-level matrix,
the squared interband matrix element of $g\cdot\sigma$ is
$|g|^2-(g\cdot d)^2/|d|^2$. Here
$|g|^2=4t^2\sin^2(k/2)$ and $g\cdot d=-2t^2\sin k$, so direct simplification
gives
\begin{equation}\label{eq:floquet-weight}
 w(k):=|\langle k,+|i[K,X]|k,-\rangle|^2
 =\frac{4a^2t^2\sin^2(k/2)}{R(k)^2}>0
 \quad(0<k<\pi).
\end{equation}
Consequently the off-diagonal part of the ccc measure is, up to the harmless
positive trace normalization,
\begin{equation}\label{eq:floquet-ccc}
 \int_{-\pi}^{\pi}w(k)
 \left[\delta_{(\lambda_+(k),\lambda_-(k))}
       +\delta_{(\lambda_-(k),\lambda_+(k))}\right]\frac{dk}{2\pi}.
\end{equation}
It is already clear that this is positive measure on one-dimensional curves.

To match the source's one-sided rectangle convention, observe that
\begin{equation}\label{eq:band-derivatives}
 \lambda_+'(k)=-\sin k\left(2s+\frac{t^2}{R(k)}\right),\qquad
 \lambda_-'(k)=-\sin k\left(2s-\frac{t^2}{R(k)}\right).
\end{equation}
Both derivatives are negative on $(0,\pi)$ because $R(k)\ge a$ and
$2s>t^2/a$. Fix $k_0=\pi/2$ and put
$(E_1^0,E_2^0)=(\lambda_+(k_0),\lambda_-(k_0))$; this point is off the
diagonal. By continuity and \eqref{eq:band-derivatives}, there are constants
$c,C,\eta>0$ such that, for $0<h<\eta$ and both signs,
\[
 ch\le\lambda_\pm(k_0-h)-\lambda_\pm(k_0)\le Ch.
\]
Hence the upper-right $\varepsilon$-square based at
$(E_1^0,E_2^0)$ contains the first curve in \eqref{eq:floquet-ccc} for every
$k=k_0-h$ with $0\le h\le\varepsilon/C$. Equation
\eqref{eq:floquet-weight} is bounded below there by a positive constant, so
the square has mass at least $c_0\varepsilon$. Division by
$\varepsilon^2$ proves the claim.
\end{proof}

\section{A completely localized Wegner example}

Work on scalar $\ell^2(\Z)$ and partition it into cells
$C_n=\{2n,2n+1\}$. Let $t>0$, let $(\omega_n)_{n\in\Z}$ be i.i.d. uniform on
$[-1,1]$, and define
\begin{align}
 (H_\omega\psi)(2n)&=\omega_n\psi(2n)+t\psi(2n+1),\label{eq:model-a}\\
 (H_\omega\psi)(2n+1)&=t\psi(2n)+\omega_n\psi(2n+1).\label{eq:model-b}
\end{align}
Thus $H_\omega$ is a bounded, self-adjoint, finite-range random Jacobi
operator with deterministic hopping. Let $X\delta_j=j\delta_j$.

Let $\tau_a$ shift the cell variables and let $U_a\delta_j=\delta_{j+2a}$.
Then $a\mapsto U_a$ is a unitary representation of $\Z$ and
$U_aH_\omega U_{-a}=H_{\tau_a\omega}$ (up to the harmless choice of shift
sign). The Bernoulli shift is ergodic. This is exactly a two-site fundamental
cell version of the abstract covariance setup in Section~2 of~\cite{cgh},
which assumes an arbitrary unitary representation of $\Z^d$. We normalize
trace per physical site by
\begin{equation}\label{eq:trace}
 \mathcal T(A)=\frac12\E\operatorname{Tr}(P_0A_\omega P_0),
\end{equation}
where $P_0$ projects onto $C_0$.

\begin{theorem}\label{thm:main}
For the ergodic family \eqref{eq:model-a}--\eqref{eq:model-b}:
\begin{enumerate}
\item every spectral energy is in the complete-localization regime;
\item the Wegner estimate $\mathcal T(E_H(J))\le |J|/2$ holds for every Borel
interval $J$;
\item the ccc measure is
\begin{equation}\label{eq:ccc-formula}
 M(dE_1,dE_2)=\frac{t^2}{4}\int_{-1}^{1}
 \left[\delta_{(u+t,u-t)}+\delta_{(u-t,u+t)}\right],du;
\end{equation}
in particular it is singular with respect to two-dimensional Lebesgue measure
and has no finite density at any interior point of either off-diagonal support
line.
\end{enumerate}
More explicitly, for $0<\varepsilon<\min\{1,2t\}$,
\begin{equation}\label{eq:blowup}
 \frac{M([t,t+\varepsilon],[-t,-t+\varepsilon])}{\varepsilon^2}
 =\frac{t^2}{4\varepsilon}\longrightarrow+\infty.
\end{equation}
\end{theorem}

\begin{proof}
On $C_n$, the Hamiltonian and velocity are
\[
 H_{\omega,n}=\begin{pmatrix}\omega_n&t\\t&\omega_n\end{pmatrix},
 \qquad
 v_n=i[H_{\omega,n},X]
   =\begin{pmatrix}0&it\\-it&0\end{pmatrix}.
\]
The velocity is independent of $\omega$. The normalized eigenvectors
\[
 |+\rangle=2^{-1/2}(1,1),\qquad
 |-\rangle=2^{-1/2}(1,-1)
\]
have respective energies $\omega_n+t$ and $\omega_n-t$. Direct calculation
gives
\begin{equation}\label{eq:weights}
 \langle+|v_n|+\rangle=\langle-|v_n|-\rangle=0,
 \qquad
 |\langle+|v_n|-\rangle|^2
 =|\langle-|v_n|+\rangle|^2=t^2.
\end{equation}
Substitution of \eqref{eq:weights} into \eqref{eq:source-ccc}, followed by the
cell expectation and the factor $1/2$ in \eqref{eq:trace}, gives
\eqref{eq:ccc-formula}. Its support is contained in the two lines
$E_1-E_2=\pm2t$, both disjoint from the diagonal.

For the rectangle in \eqref{eq:blowup}, the first branch of
\eqref{eq:ccc-formula} contributes precisely when $u\in[0,\varepsilon]$.
The transposed branch contributes nothing when $\varepsilon<2t$. Hence its
mass is $(t^2/4)\varepsilon$, proving \eqref{eq:blowup}. This also directly
precludes a finite point density in the one-sided rectangle convention used
by the source; the line support independently proves failure of absolute
continuity.

It remains to verify the two regularity properties. Every eigenfunction is
supported in one cell. Consequently $f(H_\omega)$ is cell-block-diagonal for
every bounded Borel $f$, and $P_n f(H_\omega)P_m=0$ for $n\ne m$. The
strong-dynamical-localization estimate of the source therefore holds with
finite constants at all energies (indeed, its left side vanishes between
distinct cells).

Finally, each cell has the two eigenvalues $\omega_n\pm t$. Thus, for every
interval $J$,
\[
 \E\operatorname{Tr}\bigl(P_0E_H(J)P_0\bigr)
 =\frac12|([-1,1]+t)\cap J|+\frac12|([-1,1]-t)\cap J|
 \le |J|.
\]
Dividing by two as in \eqref{eq:trace} yields
$\mathcal T(E_H(J))\le |J|/2$. This proves all assertions.
\end{proof}

\section{Verification, novelty, and limitations}

The accompanying standard-library script checks the block eigenpairs,
velocity weights, five-cell support, four rectangle scales, the uniform-law
Wegner constant, and five Floquet momenta for the signs and weight in
\eqref{eq:floquet-weight}--\eqref{eq:band-derivatives}. It reports
\texttt{PASS}. These checks are not used as proof; all identities above are
exact.

A bounded arXiv/web search on 13 August 2026 used the exact source title and
the phrases ``current--current correlation measure off-diagonal density,''
``existence of a density current--current correlation off-diagonal,'' and
``current--current correlation measure random dimer.'' It recovered the source
and Bellissard--Hislop's analytic-density theorem~\cite{bh}, but no occurrence
of this cell construction or a general negative answer. The novelty assessment
is therefore \emph{apparently new within the bounded search}, not an exhaustive
bibliographic guarantee.

Theorem~\ref{thm:periodic} answers the question negatively at the breadth of
the source's stated scalar ergodic framework, with literal one-site
translations. Theorem~\ref{thm:main} additionally shows that complete
localization plus Wegner control is insufficient in a standard unit-cell
extension: a genuinely two-energy nondegeneracy assumption is needed. Neither
theorem decides whether an off-diagonal density exists for the canonical
Anderson Hamiltonian with independent randomness at each physical site and
uniform nonzero hopping between all neighboring sites. Human review should
focus on the Floquet current calculation and this scope distinction.

\begin{thebibliography}{9}
\bibitem{cgh}
J.-M. Combes, F. Germinet, and P. D. Hislop,
\emph{Conductivity and the current--current correlation measure},
J. Phys. A: Math. Theor. \textbf{43} (2010), 474010;
\href{https://arxiv.org/abs/1008.5344}{arXiv:1008.5344}.

\bibitem{bh}
J. V. Bellissard and P. D. Hislop,
\emph{Smoothness of correlations in the Anderson model at strong disorder},
Ann. Henri Poincar\'e \textbf{8} (2007), 1--26;
\href{https://arxiv.org/abs/math-ph/0605042}{arXiv:math-ph/0605042}.
\end{thebibliography}

\end{document}
